\documentclass[12pt]{article}
\usepackage[utf8]{inputenc}
\usepackage[T1]{fontenc}
\usepackage{hyperref}
\title{Buyang, Baha Language Resources for K-12 Curriculum Development}
\author{Compiled by Claude Sonnet 5.5 (claude-sonnet-5-5) for Andr\'as Kornai}
\date{}
\begin{document}
\maketitle

\section{Introduction}
Baha Buyang (ISO 639-3: yha; Glottocode baha1256) is a Kra-Dai language of the Kra branch. Search results place it in Guangnan County, Yunnan, China (Yangliancun in Diyu Township and Anshecun in Bada Township). Speaker counts differ: about 600 in ASJP and 700 in Joshua Project. The search found no school materials, and very little else. Classification varies between sources: Glottolog places it in Southern Kra, and Ostapirat treats Paha as Central Kra.

\section{Primary Educational Resources}
No primers, textbooks or school programmes were found in search results.

\section{Linguistic and Reference Materials}
\begin{itemize}
\item Li Jinfang (1996), ``The sound system of Baha Buyang'', is cited as the main phonological description. The paper itself was not found. The entry in the ThoT tone database is: \url{https://thot.huma-num.fr/db/baha1256/tonal_status}
\item Wikipedia on the Buyang language: \url{https://en.wikipedia.org/wiki/Buyang_language}
\item Wikipedia on the Kra languages, citing Ostapirat's work on the subgroup: \url{https://en.wikipedia.org/wiki/Kra_languages}
\item ASJP word list entry for Baha Buyang: \url{https://asjp.clld.org/languages/BAHA_BUYANG}
\end{itemize}

\section{Religious and Community Resources}
\begin{itemize}
\item Joshua Project language page, which reports that Bible translation is still needed: \url{https://joshuaproject.net/index.php/languages/yha}
\end{itemize}

\section{Digital and Audio Resources}
No audio, video or corpus resources specific to Baha Buyang were found.

\section{Curriculum Development Context}
The language is endangered and is not taught in schools. Young people do not all speak it. Any curriculum would be built from scratch, starting with an orthography and a lexicon prepared with community speakers.
\end{document}